\scanpage{17}
and assumes that $i_0(-1)$ operates on $V$ as \inserted{homothety} $(-1)^n$. (It is not hard to describe $\unsure{\gamma|\mathscr{T}}$ by universal propert\deleted{\unsure{ies|y}}y as \deleted{\unsure{being}} representing some functor \inserted{$F(X)$, $X$ \deleted{\unsure{al|an}}complex analytic,} expressing that every $V$ \inserted{as above} should define, in an additive and multiplicative way for varying $V$, a Hodge-structure \inserted{$\unsure{\mathscr{V}|V}$} over $X$). If $\Gamma\subset G(\deleted{\unsure{\mathbb{Q}|\mathbb{R}}}\unsure{\mathbb{R}|\mathbb{Q}})$ is the group leaving $V_0$ fixed, then $\Gamma$ operates on $(\mathscr{V},\gamma)$.
% margin note written sideways
\margin{Note that in general $\gamma$ is \uline{not} \unsure{riemanian|riemannian} symm., but is a fiber space over a R.S. space, with fibers which are projective complex \unsure{homogeneous|homogenous} spaces under $K_{\mathbb{C}}$. This implies at least that $\gamma$ is \uline{simply} \uline{connected}, hence is the \unsure{univ.|univ-} covering of $M=\gamma/\Gamma$ if $\Gamma$ operates freely}

The modular varieties I expect to carry algebraic structure\deleted{\unsure{s}}\unsure{{}|s} \uline{over} $\mathbb{Q}$ (and morally, $\mathbb{Z}$) are the varieties $\boxed{\gamma/\Gamma}$, where $\Gamma\subset G(\mathbb{Q})\cap G(\mathbb{R}^{\unsure{0|\circ}})$ \deleted{\unsure{is}} is commensurable with $G(\mathbb{Z})$ [for some matrix representation of $G$, \deleted{\unsure{as|so}} giving a meaning to $G(\mathbb{Z})$]. However, it is possible that one will have to impose on the data some extra conditions, either arithmetic (involving existence of ``Frobenius elements'' in $\deleted{\unsure{G(\mathbb{Q})|G}}\in\unsure{.|{}}\;I(\mathbb{Q})$, $I=G/\text{\unsure{élat.|élém.} \unsure{aut.|ant.}}$),\unsure{.|{}} or geometric \inserted{on $G_{\mathbb{R}}$} (such as: $G_{\mathbb{R}}$ has \uline{no compact factor}). A typical case which fits in the general \deleted{\unsure{picture|pictures}} \inserted{description}, but for which I have no evidence so far if $\gamma/\Gamma$ has
% archive sheet number: 16
